% Commenti speciali

	% !TEX encoding = UTF-8 Unicode
	% !TEX TS-program = pdflatex
% !TeX spellcheck = it_IT
	% !TEX root = Dispensa/afed.tex

% Capitolo 2

\chapter*{Preliminari}
\addcontentsline{toc}{chapter}{Preliminari}
\markboth{PRELIMINARI}{}
\section{Spazi riflessivi}
\begin{notazione}
	Dato $X$ uno spazio normato su $\K$, indichiamo con $X' = \mathcal{L}(X;\K)$ il suo duale topologico\index[simboli]{$X'$}\index[analitico]{duale topologico} e con $X'' = \mathcal{L}(X';\K)$ il suo biduale topologico\index[simboli]{$X''$}\index[analitico]{biduale topologico}.
\end{notazione}
\begin{definizione}
	Sia $X$ uno spazio normato su $\K$. Chiamiamo \emph{iniezione canonica}\index[simboli]{$J$}\index[analitico]{iniezione canonica} di $X$ in $X''$, l'applicazione $J: X \to X''$ tale che per ogni $x \in X$, per ogni $\varphi \in X'$
	\[ \langle Jx,\varphi\rangle = \langle\varphi, x\rangle. \]
\end{definizione}
\begin{proposizione}
	Sia $X$ uno spazio normato su $\K$ e $J$ l'iniezione canonica di $X$ in $X''$. Valgono i seguenti fatti:
	\begin{enumerate}[$(a)$]
		\item $J$ è ben definita,
		\item $J$ è un'isometria, ovvero per ogni $x \in X$
		\[ \norm{Jx}_{X''} = \norm{x}_{X}, \]
		\item $J$ è lineare e iniettiva\footnote{Il nome è quindi giustificato.}.
	\end{enumerate}
\end{proposizione}
\begin{proof}
	Omettiamo la dimostrazione.
\end{proof}
\begin{definizione}
	Sia $X$ uno spazio normato su $\K$. Diciamo che $X$ è \emph{riflessivo}\index[analitico]{spazio!riflessivo} se $J$ è suriettiva.
\end{definizione}
\begin{proposizione}
	Sia $X$ uno spazio normato su $\K$ di dimensione finita. Allora $J$ è biettiva. In particolare $X$ è riflessivo.
\end{proposizione}
\begin{proof}
	È sufficiente combinare il fatto che $J$ sia iniettiva con il Teorema di Nullità + Rango.
\end{proof}
\begin{teorema}
	Sia $X$ uno spazio normato su $\K$. Allora $X$ è riflessivo se e solo se $X'$ è riflessivo.
\end{teorema}
\begin{proof}
	Omettiamo la dimostrazione.
\end{proof}
\begin{teorema}
	Siano $X$ uno spazio di Banach su $\K$ riflessivo e $Y \le X$ chiuso. Allora $Y$ è riflessivo.
\end{teorema}
\begin{proof}
	Omettiamo la dimostrazione.
\end{proof}
\begin{proposizione}\label{normrifban}
	Sia $X$ uno spazio normato su $\K$ riflessivo. Allora $X$ è di Banach.
\end{proposizione}
\begin{proof}
	Sia $(x_{h})$ di Cauchy in $X$. Allora per ogni $\varepsilon > 0$ esiste $\bar{h} \in \N$ tale che per ogni $h,k \ge \bar{h}$, $\norm{x_{h}- x_{k}}_{X}<\varepsilon$. Detta $J$ l'iniezione canonica associata a $X$,
	\[ \norm{Jx_{h} - Jx_{k}}_{X''} = \norm{J(x_{h} - x_{k})}_{X''} = \norm{x_{h} - x_{k}}_{X} < \varepsilon \]
	ossia $(Jx_{h})$ è di Cauchy in $X''$, che è completo essendo $\K$ completo. Di conseguenza esiste $T \in X''$ tale che $Jx_{h} \to T$. 
	
	Poiché $X$ è riflessivo, esiste $x \in X $ tale che $T = Jx$. Essendo $J$ un'isometria,
	\[ \norm{Jx_{h} - Jx}_{X''} = \norm{x_{h} - x}_{X}\]
	e passando al limite otteniamo che $x_{h} \to x$.
\end{proof}
\begin{esempio}
	Sia 
	\[ C_{00}= \left\lbrace (x_{h}) \in l^{\infty}(\N) \;:\; \textup{esiste } m \in \N \textup{ tale che } x_{h} = 0 \textup{ per ogni } h \ge m\right\rbrace \]
	munito della norma del $\sup$. Allora $(C_{00},\norm{\cdot}_\infty)$ non è riflessivo.\index[simboli]{$C_{00}$} 
\end{esempio}
\begin{proof}
	Consideriamo la successione in $C_{00}$ definita da
	\[ x_{j} = (1,\frac{1}{2},\frac{1}{4}, \dots, \frac{1}{2^{j}}, 0, \dots, 0, \dots). \]
	Siccome $x_{j} \to (1,\frac{1}{2},\frac{1}{4}, \dots, \frac{1}{2^{j}}, \frac{1}{2^{j+1}}, \dots) \notin C_{00}$, possiamo affermare che $C_{00}$ non è chiuso. Non può quindi essere completo. La tesi discende quindi dalla Proposizione \eqref{normrifban}.
\end{proof}
\begin{definizione}
	Sia $X$ uno spazio normato su $\K$. Diciamo che $X$ è \emph{uniformemente convesso}\index[analitico]{spazio!uniformemente convesso} se per ogni $\varepsilon > 0$ esiste $\delta >0$ tale che
	\[ \forall x,y \in X, \norm{x},\norm{y} \le 1 : \norm{x-y} \ge \varepsilon \Longrightarrow \norm*{\frac{x+y}{2}}  < 1-\delta.\]
\end{definizione}
Possiamo immaginare uno spazio normato uniformemente convesso come uno spazio normato la cui palla unitaria sia in un certo senso ``tonda''. Più precisamente, si deve avere che il punto medio di segmenti con estremi sul bordo della palla unitaria sia contenuto nell'interno della palla unitaria.

Il risultato seguente collega una nozione geometrica, l'uniforme convessità, con una nozione topologica, la riflessività. Premettiamo subito che il risultato non è invertibile, i controesempi sono molteplici.
\begin{teorema}\label{uniconvrif}
	Sia $X$ uno spazio di Banach su $\K$ uniformemente convesso. Allora $X$ è riflessivo.
\end{teorema}
\begin{proof}
	Omettiamo la dimostrazione.
\end{proof}

\section{Alcuni complementi sugli spazi di Legesgue} 
Iniziamo con alcune considerazioni di carattere generale.
\begin{proposizione}
	Consideriamo $(X,\norm{\;})$ uno spazio normato su $\K$, $(x_h)$ in $X$ e $x \in X$. I seguenti fatti sono equivalenti:
	\begin{enumerate}[$(a)$]
		\item $x_h \to x$ in $X$,
		\item per ogni sottosuccessione $(x_{h_k})$, esiste una sotto-sottosuccessione $(x_{h_{k_j}})$ tale che $x_{h_{k_j}} \to x$ in $X$.
	\end{enumerate}
\end{proposizione}
\begin{proof}\hfill
	\par\noindent
	$(a) \Longrightarrow (b)$ Evidente.
	\par\noindent
	$(b) \Longrightarrow (a)$ Supponiamo, per assurdo, che esista $\varepsilon > 0$ tale che per ogni $k > 0$ esista $h_k \ge k$ tale che $\norm{x_{h_k}-x} > \varepsilon$. Allora, per ogni sottosuccessione $(x_{h_{k_j}})$ di $(x_{h_k})$, per ogni $j \in \N$, $\norm{x_{h_{k_j}}-x} > \varepsilon$, ossia $(x_{h_{k_j}})$ non converge a $x$, assurdo.
\end{proof}

\begin{proposizione}
Siano $p \in \left[ 1,+\infty\right[ $, $k \in \N$ tale che $k \ge 2$ ed $x_1,\dots, x_k \ge 0$. Si ha
\[ \left( \sum_{i=1}^{k}x_i\right)^p \le k^{p-1}\sum_{i=1}^{k}x_i^p. \]
\end{proposizione}
\begin{proof}
Si tratta di usare la convessità di $\left\lbrace t \mapsto t^p\right\rbrace$ a più punti. Infatti, $\displaystyle\sum_{i=1}^{k}\frac{1}{k}x_i $ è una combinazione convessa di $x_1,\dots, x_k$ essendo $\displaystyle\sum_{i=1}^{k}\frac{1}{k} = 1$, 
quindi
\begin{align*}
	\left( \sum_{i=1}^{k}x_i\right)^p = \left( k\sum_{i=1}^{k}\frac{1}{k}x_i\right)^p = k^p  \left(\sum_{i=1}^{k}\frac{1}{k}x_i\right)^p \le k^p \sum_{i=1}^{k}\frac{1}{k}x_i^p = k^{p-1}\sum_{i=1}^{k}x_i^p.\qedhere
\end{align*}
\end{proof}

Nel seguito $U$ indicherà un generico sottoinsieme aperto di $\R^n$. 

\begin{teorema}\emph{\textbf{(disuguaglianza di interpolazione)}}\index[analitico]{disuguaglianza! di interpolazione}
Siano $1 \le q \le p \le r \le \infty$ e $\vartheta \in \left] 0,1\right[ $ tali che 
\[ \frac{1}{p} = \frac{\vartheta}{q} + \frac{1-\vartheta}{r}.  \]
Se $u \in L^q(U)\cap L^r(U)$, allora $u \in L^p(U)$ e 
\[ \norm{u}_{L^p(U)} \le \norm{u}_{L^q(U)}^\vartheta\norm{u}_{L^r(U)}^{1-\vartheta}.\]

\begin{proof}
Vediamo solo il caso finito. Innanzitutto,
\[ \int_{U} \abs{u}^pd\mathcal{L}^n = \int_{U} \abs{u}^{\vartheta p}\abs{u}^{(1-\vartheta)p}d\mathcal{L}^n  \]
e per la disuguaglianza di Holder\footnote{Essendo \[ 1 = \frac{\vartheta p}{q} + \frac{(1-\vartheta)p}{r}.  \]}
\[ \int_{U} \abs{u}^pd\mathcal{L}^n \le \left( \int_{U}\abs{u}^{q}d\mathcal{L}^n\right) ^{\frac{\vartheta p}{q}} \left( \int_{U}\abs{u}^rd\mathcal{L}^n\right)^{\frac{(1 -\vartheta) p}{r}} = \norm{u}_{L^q(U)}^{\vartheta p}  \norm{u}_{L^r(U)}^{(1- \vartheta) p} \]
da cui la tesi.
\end{proof}
\end{teorema}
\begin{teorema}\textbf{\emph{(disuguaglianza di Holder generalizzata)}}\footnote{Ma è strabiliante! Si potrebbe quasi dire che è \emph{over 9000}. A parte le citazioni di DragonBall (e parodie), si tratta della versione più generale della disuguaglianza di Holder valida per gli spazi normati $\left( L^p(U),\norm{\ }_p\right) $.}\index[analitico]{disuguaglianza! di Holder generalizzata}\label{holdergeneralizzato}
Consideriamo $1\le p_1,\dots, p_N \le \infty$. Sia $1\le r\le \infty$  tale che 
\[ \frac{1}{r} = \frac{1}{p_1} + \dots + \frac{1}{p_N}.\footnote{Di fatto $r$ è completamente determinato dai $p_1,\dots, p_N$.}\]
Se $f_1 \in L^{p_1}(U), \dots, f_N \in L^{p_N}(U)$, allora $f = f_1 \dots f_N \in L^{r}(U)$ e 
\[ \norm{f}_r \le \prod_{i=1}^{N} \norm{f_i}_{p_i}. \] 
\begin{proof}
	Supponiamo che per ogni $k \in \N$ tale che $k < N$ valga la tesi. Mostriamo che vale anche per $N$. Innanzitutto, a meno di riordinare gli indici ed il prodotto delle $f_i$, possiamo supporre $p_1 \le \dots \le p_N$. Se per ogni $i = 1,\dots, N$, $p_i = \infty$, allora anche $r = \infty$ e la tesi è evidente. Se invece esiste $ j \in \N$ tale che $0 <j < N$ e $p_i = \infty$ per $i = j+1, \dots, N$ e $p_i < \infty$ per $i= 1,\dots, j$, allora
	\[ \frac{1}{r} = \frac{1}{p_1} +\dots + \frac{1}{p_{j}} \] 
	ed essendo $j < N$,
	\[ \int_{U} \abs{f}^{r} d\mathcal{L}^n \le \prod_{i=j+1}^{N} \norm{f_i}_\infty^r \int_{U} \prod_{i=1}^{j}\abs{f_i}^r d\mathcal{L}^n \le \prod_{i=j+1}^{N} \norm{f_i}_\infty^r \prod_{i=1}^{j} \norm{f_i}^r_{p_i} = \left( \prod_{i=1}^N \norm{f_i}_{p_i}\right)^r \] 
	da cui 
	\[ \norm{f}_r \le \prod_{i=1}^{N}\norm{f_i}_{p_i}. \]
	Se poi $p_N < \infty$, allora per ogni $i = 1,\dots, N-1$, $p_i < \infty$ e quindi $r < \infty$. Ora, evidentemente $p_N > r$ e quindi, per la disuguaglianza di Holder\footnote{Quella classica, che applichiamo qui con esponente $p = \frac{p_N}{r}$.},
	\[ \int_{U} \abs{f}^{r} d\mathcal{L}^n = \int_{U} \abs{f_N}^r \prod_{i=1}^{N-1} \abs{f_i}^r d\mathcal{L}^n \le \left( \int_{U} \abs{f_N}^{p_N} d\mathcal{L}^n\right) ^{\frac{r}{p_N}} \left(\int_{U} \prod_{i=1}^{N-1} \abs{f_i}^{\frac{p_Nr}{p_N - r}}d\mathcal{L}^n\right) ^{\frac{p_N-r}{p_N}}. \]
	Ora, dal fatto che 
	\begin{align*}
		N-1 &< N, & \frac{p_Nr}{p_N - r} &> 1, & \frac{p_N-r}{p_N r} &= \frac{1}{p_1} + \dots +\frac{1}{p_{N-1}},
	\end{align*}
	possiamo applicare l'ipotesi induttiva forte e ottenere
	\[\int_{U} \prod_{i=1}^{N-1} \abs{f_i}^{\frac{p_Nr}{p_N - r}}d\mathcal{L}^n \le \prod_{i=1}^{N-1} \left( \int_{U} \abs{f_i}^{p_i}d\mathcal{L}^n\right) ^{\dfrac{\frac{p_Nr}{p_N -r}}{p_i}}.\]
	Pertanto,
	\[ \int_{U} \abs{f}^{r} d\mathcal{L}^n \le \left( \int_{U} \abs{f_N}^{p_N} d\mathcal{L}^n\right) ^{\frac{r}{p_N}}  \left(\prod_{i=1}^{N-1} \norm{f_i}_{p_i} \right) ^{r} = \left(\prod_{i=1}^{N} \norm{f_i}_{p_i} \right) ^{r} \]
	da cui la tesi.
\end{proof}
\end{teorema}

\begin{osservazione}
È possibile estendere il Teorema \eqref{holdergeneralizzato} anche per $r, p_1,\dots, p_N \in \left] 0,1\right[ $. Tuttavia si ricorda che in questo caso non abbiamo a che fare con delle norme in quanto non vale la disuguaglianza triangolare, pertanto il Teorema \eqref{holdergeneralizzato} è di fatto la versione più generale se si vuole lavorare con delle norme. In ogni caso la generalizzazione è diretta, lo stesso enunciato può essere scritto per $r, p_1,\dots, p_N \in \left] 0,1\right[ $ intendendo appropriatamente il simbolo $\norm{\ }$.
\end{osservazione}
Presentiamo ora un criterio di compattezza per gli spazi di Lebesgue\footnote{Così come il Teorema di Ascoli--Arzelà per le funzioni continue su un compatto. In effetti la dimostrazione richiede proprio l'uso del Teorema sopra citato.}.
\begin{notazione}\index[simboli]{$\tau_{y}(f)$}
	Per ogni $x \in \R^n$, per ogni $y \in \R^n$, indichiamo con 
	\[ \tau_{y}(f)(x) = f(x+y) \]
	la traslata di vettore $y$ di una qualsiasi applicazione $f$.
\end{notazione}

\begin{teorema}\emph{\textbf{(di Kolmogorov--M. Riesz--Fréchet)}}\index[analitico]{teorema!di Kolmogorov--M. Riesz--Fréchet}
	Siano $1 \le p < \infty$ ed $\mathcal{F}$ un sottoinsieme di $L^p(\R^{n})$. Se $\mathcal{F}$ è limitato e per ogni $\varepsilon > 0$, esiste $\delta > 0$ tale che 
	\[ \forall\; f \in \mathcal{F},\forall\; y \in \R^n: \abs{y} < \delta \Longrightarrow \norm{\tau_{y}(f)-f}_{p} < \varepsilon \]
	allora per ogni sottoinsieme misurabile $\Omega$ di $\R^n$ di misura finita, $\overline{\mathcal{F}_{| \Omega}}$ è compatto in $L^p(\Omega)$.
\end{teorema}
\begin{proof}
	Omettiamo la dimostrazione.
\end{proof}

\begin{corollario}
	Siano $1 \le p < \infty$ ed $\mathcal{F}$ un sottoinsieme di $L^p(U)$. Se $\mathcal{L}^n(U) < +\infty$, $\mathcal{F}$ è limitato e per ogni $\varepsilon > 0$, esiste $\delta > 0$ tale che 
	\[ \forall\; f \in \mathcal{F},\forall\; y \in \R^n: \abs{y} < \delta \Longrightarrow \norm{\tau_{y}(f)-f}_{p} < \varepsilon \]
	allora $\overline{\mathcal{F}}$ è compatto in $L^p(U)$.
\end{corollario}
\begin{proof}
	Si tratta di una riscrittura del Teorema di Kolmogorov.
\end{proof}
\begin{corollario}\label{corollarioKolmogorov}
	Siano $1 \le p < \infty$ ed $\mathcal{F}$ un sottoinsieme di $L^p(\R^n)$. Se $\mathcal{F}$ è limitato, per ogni $\varepsilon > 0$, esiste $\delta > 0$ tale che 
	\[ \forall\; f \in \mathcal{F},\forall\; y \in \R^n: \abs{y} < \delta \Longrightarrow \norm{\tau_{y}(f)-f}_{p} < \varepsilon \]
	ed esiste un sottoinsieme limitato $U$ di $\R^n$ misurabile tale che 
	\[ \forall\; f \in \mathcal{F}: \norm{f}_{L^p(\R^n\setminus U)} < \varepsilon \] 
	allora $\overline{\mathcal{F}}$ è compatto in $L^p(\R^n)$.
\end{corollario}
\begin{proof}
	Omettiamo la dimostrazione.
\end{proof}
\begin{osservazione}
	Il viceversa del Corollario \eqref{corollarioKolmogorov} è anch'esso vero. Pertanto abbiamo una caratterizzazione piena dei sottoinsiemi compatti di $L^p(\R^n)$.
\end{osservazione}

Passiamo quindi ad alcune considerazioni sulla convergenza.
\begin{proposizione}\label{convergenzaLq}
	Supponiamo che $\mathcal{L}^n(U)< +\infty$ e $1 < p < \infty$.  Sia $(f_{h})$ una successione limitata in $L^p(U)$. Se $f_{h} \to f$ \emph{q.o.} in $U$, allora per ogni $q < p$, $f_{h} \to f$ in $L^q(U)$.
\end{proposizione}
\begin{proof}
	Dalla limitatezza di $(f_{h})$ innanzitutto discende l'esistenza di $K > 0$ tale che per ogni $h \in \N$
	\[ \norm{f_{h}}_{p}\le K. \]
	Ora, per il Lemma di Fatou
	\[ \int_{U}\abs{f(x)}^p d\mathcal{L}^n(x) \le \underset{h}{\liminf} \int_{U} \abs{f_h(x)}^pd\mathcal{L}^n(x) = \underset{h}{\liminf} \norm{f_{h}}_{p}^p \le K^p \]
	da cui $f \in L^p(U)$. Mostriamo che per ogni $q \in \left[  1,p \right[$
	\[ \lim_{h} \int_{U}\abs{f_h(x)-f(x)}^qd\mathcal{L}^n(x) = 0. \]
	
	Fissiamo $M > 0 $ e prendiamo 
	\[ U_{h,M} = \Set{ x \in U: \abs{f_h(x) - f(x)} < M}. \]
	Ora,
	\[ \int_{U}\abs{f_h(x)-f(x)}^qd\mathcal{L}^n(x) =  \int_{U_{h,M}}\abs{f_h(x)-f(x)}^qd\mathcal{L}^n(x) + \int_{U\setminus U_{h,M}}\abs{f_h(x)-f(x)}^qd\mathcal{L}^n(x),\]
	e per il Teorema della convergenza dominata\footnote{Indichiamo con $g_h(x) = \abs{f_h(x)-f(x)}^q\chi_{U_{h,M}}$,
		\begin{align*}
			g_{h} &\to 0 \textup{ q.o. in } U & \abs{g_h} \le M^q \in L^1(U)
		\end{align*} 
		in quanto $\mathcal{L}^n(U) < +\infty$.}, a $M$ fissato,
	\[ \lim_{h} \int_{U_{h,M}}\abs{f_h(x)-f(x)}^qd\mathcal{L}^n(x) = 0.\]
	Invece, 
	\[ \begin{split}
		\int_{U\setminus U_{h,M}}\abs{f_h(x)-f(x)}^qd\mathcal{L}^n(x) &= \int_{U\setminus U_{h,M}}\abs{f_h(x)-f(x)}^p \abs{f_h(x)-f(x)}^{q-p}d\mathcal{L}^n(x) \le \\
		&\le \frac{1}{M^{p-q}}\int_{U\setminus U_{h,M}}\abs{f_h(x)-f(x)}^pd\mathcal{L}^n(x) \le \frac{C}{M^{p-q}}.
	\end{split} \]
	Allora, 
	\[ \int_{U}\abs{f_h(x)-f(x)}^qd\mathcal{L}^n(x) \le \int_{U_{h,M}}\abs{f_h(x)-f(x)}^qd\mathcal{L}^n(x) + \frac{C}{M^{p-q}} \]
	e se $h \to +\infty$
	\[ \int_{U}\abs{f_h(x)-f(x)}^qd\mathcal{L}^n(x) \le \frac{C}{M^{p-q}}. \]
	Se ora $M \to +\infty$ otteniamo la tesi.
\end{proof}

\begin{definizione}
	Siano $\mu$ una misura esterna su $\R^n$ ed $E \subseteq \R^n$ misurabile. Siano $(f_{h})$ una successione di applicazioni misurabili da $E$ in $\R$ ed $f: E \to \R$ misurabile. Diciamo che $f_h$ \emph{converge quasi uniformemente} a $f$ in $E$\index[analitico]{convergenza!quasi uniforme} se per ogni $\varepsilon>0$ esiste $E_{\varepsilon}\subseteq E$ tale che 
\begin{align*}
	\mu(E\setminus E_{\varepsilon})&< \varepsilon,  & \lim_{h}\norm{f_{h}-f}_{L^{\infty}(E_{\varepsilon},\mu)} &=0.
\end{align*} 
\end{definizione}
Chiaramente, la convergenza quasi uniforme implica la convergenza q.o. in $E$. Nel caso di misura finita, il seguente risultato realizza l'implicazione contraria.

\begin{teorema}\emph{\textbf{(di Severini--Egorov)}}\index[analitico]{teorema! di Severini--Egorov}
	Siano $\mu$ una misura esterna su $\R^n$ ed $E \subseteq \R^n$ misurabile tale che $\mu(E)< +\infty$. Sia $(f_{h})$ una successione di applicazioni misurabili da $E$ in $\R$. Se $f_{h} \to f$ \emph{q.o.} in $E$, allora $f_h \to f$ quasi uniformemente in $E$.
\end{teorema}
\begin{proof}
Sia $N \subseteq E$ tale che $\mu(N) = 0$ e $f_n \to f$ puntualmente in $E \setminus N$. Per ogni $k \in \N$ consideriamo la successione crescente $(A^{(k)}_m)$ di sottoinsiemi di $E$ tali che
\[ A^{(k)}_m = \bigcap_{n \ge m}\left\lbrace x \in E \setminus N: \abs{f_n(x)-f(x)}\le \frac{1}{k+1}\right\rbrace.  \]
Chiaramente, fissato $k \in \N$
\[ \bigcup_{m \in \N} A^{(k)}_m = E \setminus N.\]
Siccome siamo in misura finita, a $k \in \N$ fissato, $\mu(A^{(k)}_m) \to \mu(E)$, per ogni $\varepsilon >0$ esiste $m_k$ tale che $\mu(E \setminus  A^{(k)}_{m_k}) \le \frac{\varepsilon}{2^{k+2}}$ e 
\[ \forall \; x \in  A^{(k)}_{m_k}, \forall \; n \ge m_k: \abs{f_n(x)-f(x)} \le \frac{1}{k+1}.\]
Posto
\[ E_\varepsilon = \bigcap_{k \in \N} A^{(k)}_{m_k}, \]
allora 
\[ \lim_{h}\norm{f_{h}-f}_{L^{\infty}(E_{\varepsilon},\mu)} =0 \]
e 
\[ \mu(E \setminus E_\varepsilon) \le \sum_{k=0}^{\infty}\frac{\varepsilon}{2^{k+2}} = \frac{\varepsilon}{2} <\varepsilon,\]
da cui la tesi.
\end{proof}
In generale, se togliamo l'ipotesi di misura finita, tuttavia, il risultato non è più vero.
\begin{esempio}
	Sia $E = \R^+$. Allora $\mathcal{L}^1(E) = +\infty$. Consideriamo 
	\[ f_h(x) = \chi_{\left] h,h+1\right[ }(x). \]
	Evidentemente $f_h(x) \to 0$. Supponiamo per assurdo valga il Teorema di Severini--Egorov, allora per ogni $\varepsilon > 0$, esiste $E_{\varepsilon} \subseteq \R^+$ tale che $\mathcal{L}^1(\R^+\setminus E_{\varepsilon}) < \varepsilon$ e 
	\[ \lim_{h} \norm{f_h}_{L^{\infty}(E_{\varepsilon})} = 0. \]
	Allora, definitivamente in $h$, 
	\[ E_{\varepsilon} \cap \left] h,h+1\right[ = \emptyset, \]
	quindi $E_{\varepsilon}$ limitato. Allora $\mathcal{L}^1(\R^+\setminus E_{\varepsilon}) = +\infty$, assurdo.
\end{esempio}


\begin{osservazione}
Mediante il Teorema di Severini--Egorov, possiamo costruire una diversa dimostrazione della Proposizione \eqref{convergenzaLq}. 

Per il Teorema di Severini--Egorov, per ogni $\varepsilon > 0$, esiste $U_{\varepsilon}$ tale che $\mathcal{L}^n(U\setminus U_{\varepsilon})<\varepsilon$ e 
\[ \lim_{h}\norm{f_{h}-f}_{L^{\infty}(U_{\varepsilon})} =0. \]
Ora, per la disuguaglianza di Holder,
\[\begin{split}
	 \int_{U}\abs{f_h(x)-f(x)}^qd\mathcal{L}^n(x) &=   \int_{U_{\varepsilon}}\abs{f_h(x)-f(x)}^qd\mathcal{L}^n(x) + \int_{U \setminus U_{\varepsilon}}\abs{f_h(x)-f(x)}^qd\mathcal{L}^n(x) \le \\
	 &\le \norm{f_h-f}_{L^\infty(U_{\varepsilon})}^q\mathcal{L}^n(U_{\varepsilon}) + \\
	 & + \mathcal{L}^n(U\setminus U_{\varepsilon}))^{\frac{p-q}{p}}\left( \int_{U \setminus U_{\varepsilon}} \abs{f_h(x)-f(x)}^pd\mathcal{L}^n(x)\right) ^{\frac{q}{p}} < \\
	 & < \norm{f_h-f}_{L^\infty(U_{\varepsilon})}^q\mathcal{L}^n(U) + \widehat{C}\varepsilon^{\frac{p-q}{p}}
\end{split} \]
Ora, passando prima al limite per $h \to +\infty$ e poi per $\varepsilon \to 0$ si ottiene la tesi.
\end{osservazione}

\begin{definizione}\index[analitico]{convergenza! in misura}
	Siano $\mu$ una misura esterna su $\R^n$ ed $E \subseteq \R^n$ misurabile. Siano $(f_{h})$ una successione di applicazioni misurabili da $E$ in $\R$ ed $f: E \to \R$ misurabile. Diciamo che $f_h$ \emph{converge in misura} a $f$, e lo indichiamo con $f_h \overset{\mu}{\to} f$\index[simboli]{$\overset{\mu}{\to}$} se, per ogni $\varepsilon > 0 $, detto 
\[ E_{\varepsilon} = \left\lbrace  x \in U: \abs{f_h(x) - f(x)} > \varepsilon\right\rbrace, \]
risulta
\[ \lim_{h} \mu(E_{\varepsilon}) = 0. \]
\end{definizione}
\begin{proposizione}
	Siano $\mu$ una misura esterna su $\R^n$ ed $E \subseteq \R^n$ misurabile. Siano $(f_{h})$ una successione di applicazioni misurabili da $E$ in $\R$ ed $f: E \to \R$ misurabile. Se $f_h \overset{\mu}{\to} f$, allora esiste $(f_{h_k})$ tale che $f_{h_k}\to f$ \emph{q.o.} in $E$.
\end{proposizione}
\begin{proof}
Per ogni $k \in \N$, esiste $n_k$ tale che 
\[ \mu\left( \left\lbrace x \in E: \abs{f_n(x)-f(x)}>\frac{1}{k}\right\rbrace \right) \le \frac{1}{k^2}. \]
Siano
\[ A_k = \left\lbrace x \in E: \abs{f_n(x)-f(x)}>\frac{1}{k}\right\rbrace   \]
e 
\[ A = \bigcap_{m \in \N} \bigcup_{k \ge m} A_k. \]
La tesi segue dal fatto che $\mu(A) \le \displaystyle \sum_{k=m}^{\infty} \frac{1}{k^2} \to 0$ e $f_{n_k}(x) \to f(x)$ in  $E\setminus A$ in quanto 
\[ E \setminus A = \bigcup_{m \in N} \bigcap_{k \ge m} \left\lbrace x \in E: \abs{f_{n_k}(x)-f(x)} \le \frac{1}{k} \right\rbrace. \qedhere\]
\end{proof}

\begin{proposizione}
	Siano $\mu$ una misura esterna su $\R^n$ ed $E \subseteq \R^n$ misurabile tale che $\mathcal{L}^n(E)<+\infty$. Siano $(f_{h})$ una successione di applicazioni misurabili da $E$ in $\R$ ed $f: E \to \R$ misurabile. Se $f_h \to f$ \emph{q.o.} in $E$, allora $f_h \overset{\mu}{\to} f$.
\end{proposizione}
\begin{proof}
Fissiamo $\varepsilon >0$. Per il Teorema di Severini--Egorov, per ogni $\delta >0$ esiste $E_{\delta}\subseteq E$ tale che 
\begin{align*}
	\mu(E\setminus E_{\delta})&<\delta,  & \lim_{h}\norm{f_{h}-f}_{L^{\infty}(E_{\delta},\mu)} &=0.
\end{align*} 
In particolare, esiste $n_\varepsilon$ tale che per ogni $h \ge n_\varepsilon$ e per ogni $x \in E_\delta$
\[ \abs{f_h(x)-f(x)}< \varepsilon. \]
Allora, per $h \ge n_\varepsilon$
\begin{align*}
 \mu\left( \left\lbrace x \in E: \abs{f_h(x)-f(x)}>\varepsilon\right\rbrace \right) &\le \mu\left( \left\lbrace x \in E_\delta: \abs{f_h(x)-f(x)}>\varepsilon \right\rbrace \right) + \\
 &+\mu\left( \left\lbrace x \in E\setminus E_\delta: \abs{f_h(x)-f(x)}>\varepsilon \right\rbrace \right) \le \\
 &\le \mu(\emptyset) + \mu(E \setminus E_\delta) < \delta.\qedhere
\end{align*}
\end{proof}

\begin{proposizione}
Siano $1\le p < \infty$ e $(f_h)$ una successione in $L^p(U,\mu)$. Se 
\[ \lim_{h}\norm{f_h-f}_{p} = 0 \]
allora $f_h \overset{\mu}{\to} f$.
\end{proposizione}
\begin{proof}
Siccome
	\[ \int_{U} \abs{f_h - f}^p d\mu \ge \int_{E_{\varepsilon}}\abs{f_h - f}^p d\mu > \varepsilon^p\mu(E_{\varepsilon}),\]
fissato $\varepsilon$, se $h \to +\infty$ risulta che $\mu(E_{\varepsilon}) \to 0$.
\end{proof}
\begin{esempio}
Sia $1<p<\infty$. Consideriamo la successione definita da $f_h: \R^+ \to \R$ tale che
\[ f_h(x) = h\chi_{\left] 0,\frac{1}{h}\right[ }(x). \]
Evidentemente $f_h \to 0$ \emph{q.o.} ma
\[ \int_{\R^+}\abs{f_h}^pd\mathcal{L}^1 = h^{p-1} \]
per cui $f_h$ non converge in $L^p(\R^+)$. Tuttavia
\[ \mathcal{L}^1 \left( \left\lbrace x \in \R^+: f_h(x) > \varepsilon\right\rbrace \right)  = \mathcal{L}^1\left( \left\lbrace x \in \R^+: h\chi_{\left] 0,\frac{1}{h}\right[ }(x) > \varepsilon\right\rbrace \right)  \le \mathcal{L}^1\left( \left] 0,\frac{1}{h}\right[\right) = \frac{1}{h}\]
per cui  $f_h \overset{\mathcal{L}^1}{\to} 0$.
\end{esempio}

Vediamo ora di completare la panoramica sulle proprietà degli spazi di Lebesgue.
\begin{teorema}\emph{\textbf{(Disuguaglianze di Clarkson)}}\index[analitico]{disuguaglianza!di Clarkson} Sia $p \in \left[ 1, \infty\right[ $. Siano $E \subseteq \R^{n} $ misurabile ed $f,g \in L^{p}(E)$. Valgono i seguenti fatti:
	\begin{enumerate}[$(a)$]
		\item se $p \ge 2$, 
		\[ \norm*{\frac{f+g}{2}}^{p}_{p} + \norm*{\frac{f-g}{2}}^{p}_{p} \le \frac{1}{2}\left( \norm{f}^{p}_{p} +  \norm{g}^{p}_{p}\right) \qquad \textup{(Disuguaglianza di Clarkson I)}\]
		\item se $1 \le p < 2$, 
		\[ \norm*{\frac{f+g}{2}}^{p'}_{p} + \norm*{\frac{f-g}{2}}^{p'}_{p} \le \left( \frac{1}{2}\norm{f}^{p}_{p} +  \frac{1}{2}\norm{g}^{p}_{p}\right)^{\frac{1}{p-1}} \quad \textup{(Disuguaglianza di Clarkson II)}\]
	\end{enumerate}
\end{teorema}
\begin{proof} \hfill
	\par\noindent
	$(a)$ Mostriamo innanzitutto che per ogni $ x,y \ge 0$
	\begin{equation}\label{eq1}
		x^{p} + y^{p} \le (x^{2} + y^2)^{\frac{p}{2}}.
	\end{equation} 
		Se $y = 0$, la disuguaglianza è evidentemente vera. Supponiamo quindi $y \ne 0$ e studiamo la veridicità di 
	\[ \left(\frac{x}{y}\right)^p + 1 \le \left( \left( \frac{x}{y}\right) ^2 + 1\right) ^{\frac{p}{2}}  \]
	o, equivalentemente (introducendo $z = \frac{x}{y}$), di 
	\[ z^p +1 \le (z^2 + 1)^{\frac{p}{2}}. \]
	Posto $\varphi(z) = (z^2 + 1)^{\frac{p}{2}} - z^p - 1 $, è sufficiente mostrare che $\varphi(z) \ge 0$ per $z \ge 0$. Grazie al fatto che $p \ge 2$, 
	\[ \varphi'(z) = p\left( z(z^2 + 1)^{\frac{p-2}{2}} - z^{p-1}\right) \ge 0. \]
	La funzione $\varphi$ è dunque crescente e, dato che $\varphi(0) = 0$, ne deduciamo che $\varphi(z) \ge 0$ per ogni $z \ge 0$.
	
	Siano ora $x = \abs*{\frac{a+b}{2}}, y = \abs*{\frac{a-b}{2}}$, con $a,b \in \R$. Allora applicando \eqref{eq1}, otteniamo 
	\[ \abs*{\frac{a+b}{2}}^{p} + \abs*{\frac{a-b}{2}}^{p} \le \left( \abs*{\frac{a+b}{2}}^{2} + \abs*{\frac{a-b}{2}}^2\right) ^{\frac{p}{2}} = \left( \frac{1}{2}a^{2} + \frac{1}{2}b^{2}\right)^{\frac{p}{2}}\le \frac{1}{2} \left( a^{p} + b^{p}\right), \]
	dove l'ultima maggiorazione è giustificata dalla convessità dell'applicazione $t \mapsto t^{\frac{p}{2}}$ su $t\ge 0$, essendo $p\ge2$.
	
	Non resta che porre $a = f(x)$ e $b= g(x)$ e integrare ambo i membri per ottenere la disuguaglianza cercata.
	\par\noindent
	$(b)$ Omettiamo la dimostrazione.\qedhere
\end{proof}
\begin{teorema}
	Siano $p \in \left[ 2, \infty \right[ $ ed $E \subseteq \R^{n} $ misurabile. Allora $L^{p}(E)$ è riflessivo.
\end{teorema}
\begin{proof}
	Mostriamo che $L^{p}(E)$ è uniformemente convesso; fatto questo, la tesi discende da \eqref{uniconvrif}. Siano $\varepsilon > 0$ e $f,g \in L^{p}(E)$ tali che $\norm{f}_p,\norm{g}_p \leq 1$ e $\norm{f-g}_{p} > \varepsilon$. Per la prima disuguaglianza di Clarkson,
\[ \norm*{\frac{f+g}{2}}^{p}_{p} \le \frac{1}{2}\left( \norm{f}^{p}_{p} +  \norm{g}^{p}_{p}\right) -  \frac{1}{2^{p}}\norm*{f-g}^{p}_{p} < 1 - \frac{\varepsilon^{p}}{2^{p}}. \]
Allora
\[ \norm*{\frac{f+g}{2}}_{p} < \left(1 - \frac{\varepsilon^{p}}{2^{p}}\right) ^{\frac{1}{p}} = 1 - 1 + \left(1 - \frac{\varepsilon^{p}}{2^{p}}\right) ^{\frac{1}{p}} = 1 - \left( 1 - \left(1 - \frac{\varepsilon^{p}}{2^{p}}\right) ^{\frac{1}{p}}\right).\] 
Ponendo $\delta = 1 - \left(1 - \frac{\varepsilon^{p}}{2^{p}}\right) ^{\frac{1}{p}} > 0$ e osservando che non dipende da $f,g$, si conclude la dimostrazione.
\end{proof}

	Per completare il quadro e dimostrare che anche gli spazi $L^p$ con $1 < p < 2$ sono riflessivi, operiamo ora per dualità, come mostrato nel seguente Teorema.

\begin{teorema}\label{1,2rif}
	Siano $p \in \left]  1, 2 \right[ $ ed $E \subseteq \R^{n} $ misurabile. Allora $L^{p}(E)$ è riflessivo. 
\end{teorema}

\begin{proof}
	Innanzitutto, se $p \in \left]  1, 2 \right[ $ allora $p' \in \left]  2, \infty \right[$, sicché $L^{p'}(E)$ è riflessivo. Di conseguenza, anche $\left( L^{p'}(E)\right)'$ è riflessivo per \eqref{teo:reflexive_iff}. Consideriamo poi l'applicazione lineare $T : L^{p}(E) \to \left( L^{p'}(E)\right)' $ tale che per ogni $u \in L^{p}(E) $, per ogni $f \in L^{p'}(E)$,
	\[ \langle Tu,f\rangle = \int_{E}^{} f(x)u(x)d\mathcal{L}^{n}(x).\]
	Risulta, dalla disuguaglianza di H\"older, che
	\[ \abs{\langle Tu,f\rangle} \le \norm{u}_{p}\norm{f}_{p'} \]
	da cui $\norm{Tu}_{\left( L^{p'}(E)\right)'} \le \norm{u}_{p}$. D'altro canto, se $\norm{u}_{p}=0$, allora anche $\norm{Tu}_{\left( L^{p'}(E)\right)'} = 0$; se invece $\norm{u}_{p}>0$, allora prendendo 
	\[ f_{0}(x) = \frac{\abs{u(x)}^{p-2}u(x)}{\norm{u}^{p-1}_{p}}, \]
	risulta $\abs{\langle Tu,f_{0}\rangle} = \norm{u}_{p}$, da cui segue che
	\[ \norm{Tu}_{\left( L^{p'}(E)\right)'} = \sup\left\lbrace \abs{\langle Tu,f\rangle} \; :\; f \in L^{p'}(E) \;,\; \norm{f}_{p'} \le 1\right\rbrace \ge \abs{\langle Tu,f_{0}\rangle} = \norm{u}_{p}. \]
	Se ne deduce che $\norm{Tu}_{\left( L^{p'}(E)\right)'} = \norm{u}_{p} $, ovvero $T$ è un'isometria. %In particolare, è iniettiva.
	
	Dal fatto che $T$ è un'isometria, si verifica facilmente che $T(L^{p}(E)) \le \left( L^{p'}(E)\right)'$ chiuso. Allora, per \eqref{teo:refl_subspace}, $T(L^{p}(E))$ è riflessivo. 
\end{proof}
\begin{osservazione}
	Tramite la seconda disuguaglianza di Clarkson si può ottenere una dimostrazione alternativa del Teorema precedente.
\end{osservazione}

Vediamo ora di andare a studiare la rappresentazione di forme lineari e continue di $L^p$.
\begin{lemma}\emph{\textbf{(di caratterizzazione topologica della densità)}}\index[analitico]{lemma!di caratterizzazione topologica della densità}
	Siano $X$ uno spazio di Banach su $\K$ e $Y \le X$. I seguenti fatti sono equivalenti:
	\begin{enumerate}[$(a)$]
		\item $\overline{Y} = X$, cioè $Y$ è denso in $X$, 
		\item per ogni $\varphi \in X'$
		\[  \left( \langle \varphi, y \rangle = 0 \textup{ per ogni } y \in Y\right) \Longrightarrow \varphi = 0.\]
	\end{enumerate}
\end{lemma}

\begin{proof}\hfill
	\par\noindent
	$(a) \Longrightarrow (b)$ Sia $x \in X$. Per densità esiste $(y_{j})$ in $Y$ tale che $y_{j} \to x$. Per ogni $\varphi \in X'$, per continuità,
	\[ \langle \varphi, x \rangle = \langle \varphi, \lim\limits_{j} y_{j} \rangle = \lim\limits_{j} \langle \varphi,  y_{j} \rangle =  \lim\limits_{j} 0 = 0.\]
	\par\noindent
	$(b) \Longrightarrow (a)$ Omettiamo la dimostrazione.
\end{proof}

\begin{teorema}
	Siano $p \in \left]  1, \infty \right[ $ ed $E \subseteq \R^{n} $ misurabile. Allora $\left( L^{p}(E)\right)' \cong L^{p'}(E)$.
\end{teorema}

\begin{proof}
	Consideriamo l'applicazione lineare $T : L^{p'}(E) \to \left( L^{p}(E)\right)'$ tale che per ogni $u \in L^{p'}(E)$, per ogni $f \in L^{p}(E)$
	\[ \langle Tu,f\rangle = \int_{E}^{} f(x)u(x)d\mathcal{L}^{n}(x).\]
	Con passaggi analoghi a quelli eseguiti nel Teorema \eqref{1,2rif}, possiamo verificare che $T$ è un'isometria e che $T(L^{p'}(E)) \le \left( L^{p}(E)\right)' $ chiuso. Mostriamo che $T$ è anche suriettiva.
	
	Sia $\varphi \in \left( L^{p}(E)\right)'' $ e supponiamo che per ogni $f \in L^{p'}(E)$ valga
	\[ \langle \varphi, Tf \rangle = 0. \]
	Per riflessività di $\left( L^{p}(E)\right)''$ esiste $g \in L^{p}(E)$ tale che $\varphi = Jg$ e quindi, per definizione di iniezione canonica,
	\[ \langle \varphi, Tf \rangle = \langle Jg, Tf \rangle = \langle Tf, g \rangle = \int_{E}^{}f(x)g(x)d\mathcal{L}^{n}(x) = 0 \]
	per ogni $f \in L^{p'}(E)$. In particolare 
	\[ \int_{E}^{}f(x)g(x)d\mathcal{L}^{n}(x) = 0 \]
	per ogni $f \in C^{\infty}_{c}(E)$, da cui $g = 0$ q.o. Ciò significa che anche $\varphi = 0$ e, per il Lemma di caratterizzazione topologica della densità, $T(L^{p'}(E))$ è denso in $X$. Dal fatto che $T(L^{p'}(E))$ è chiuso, si conclude la suriettività, ovvero che
	\[ T(L^{p'}(E)) = \left( L^{p}(E)\right)', \]
	da cui la tesi.
\end{proof}

\begin{teorema}\label{dualediL1}
	Sia $E \subseteq \R^{n} $ misurabile. Allora $\left( L^{1}(E)\right)' \cong L^{\infty}(E)$.
\end{teorema}

\begin{proof}
	Sia $\varphi : L^{1}(E) \to \R$ lineare e continua. Mostriamo che esiste un'unica $u \in L^{\infty}(E)$ tale che 
	\[ \langle \varphi, f \rangle = \int_{E}f(x)u(x)d\mathcal{L}^{n}(x). \]
	Sia $(E_{h})$ una successione di sottoinsiemi misurabili di $E$ tali che 
	\[ E = \bigcup_{h=1}^{\infty} E_{h}, \]
	$\mathcal{L}^{n}(E_{h}) < +\infty$ e $E_{0} = \emptyset$ e costruiamo una funzione $\vartheta \in L^{2}(E)$ tale che per ogni $h \in \N\setminus\left\lbrace 0 \right\rbrace $ esiste $\varepsilon_{h} > 0$ per cui $\vartheta|_{E_{h}} \ge \varepsilon_{h}$. Per fare ciò è sufficiente porre
	\[ \vartheta|_{E_{1}} = \alpha_{1}, \]
	\[ \vartheta|_{E_{i}\setminus E_{i-1}} = \alpha_{i} \textup{ per ogni } i >1,\]
	dove, per ogni $j \in \N\setminus\left\lbrace 0 \right\rbrace$,
	\[ \alpha_{j} \le \dfrac{\left( \mathcal{L}^{n}\left( E_{j}\setminus E_{j-1}\right)\right) ^{-\frac{1}{2}}}{j^{\beta}} \]
	con $\beta > \frac{1}{2}$, di modo che 
	\begin{equation*}
		\begin{split}
			\int_{E} \vartheta^{2}(x)d\mathcal{L}^{n}(x)&=  \int_{\bigcup_{j=1}^{\infty} E_{j}} \vartheta^{2}(x)d\mathcal{L}^{n}(x) \\ 
			&=\sum_{j = 1}^{\infty}\int_{E_{j}\setminus E_{j-1}} \vartheta^{2}(x)d\mathcal{L}^{n}(x) \\
			&=  \sum_{j = 1}^{\infty} \alpha^{2}_{j} \mathcal{L}^{n}\left( E_{j}\setminus E_{j-1}\right) \\
			&= \sum_{j=1}^\infty \frac{1}{j^{2\beta}} < +\infty.
		\end{split}
	\end{equation*}
	Consideriamo ora l'applicazione 
	\[ \begin{cases}
		L^{2}(E) \to \R \\
		f \longmapsto \langle \varphi, \vartheta f \rangle
	\end{cases}, \]
	ben definita in quanto $\vartheta, f \in L^{2}(E)$. Per il Teorema di F. Riesz esiste $v \in L^{2}(E)$ tale che per ogni $f \in L^{2}(E)$
	\begin{equation}\label{Riesz}
		\langle \varphi,\vartheta f \rangle = \int_{E}f(x)v(x)d\mathcal{L}^{n}(x).
	\end{equation}
	Sia $u = \frac{v}{\vartheta}$. Effettuando la scelta $f = \chi_{h}\frac{g}{\vartheta}$ in \eqref{Riesz} con $g \in L^{\infty}(E)$ e $\chi_{h} = \chi_{E_{h}}$, e osservando che la scelta è effettivamente ammissibile, risulta che per ogni $h \in \N$
	\begin{equation}\label{condizioneimportante}
		\forall \; g \in L^{\infty}(E): 	\langle \varphi,\chi_{h} g \rangle = \int_{E}\chi_{h}(x)g(x)u(x)d\mathcal{L}^{n}(x).
	\end{equation}
	Mostriamo che $u \in L^{\infty}(E) $. Consideriamo $C > \norm{\varphi}_{(L^{1}(E))'}$, 
	\[ A = \left\lbrace x \in E \,:\; \abs{u(x)} \ge C \right\rbrace  \]
	e, scegliendo $g = \chi_{A}\frac{u}{\abs{u}} \in L^{\infty}(E)$ in \eqref{condizioneimportante}, osserviamo che
	\[
	\begin{split}
		\langle \varphi, \chi_{h}\chi_{A}\frac{u}{\abs{u}}\rangle & = \int_{E}u(x)\chi_{h}(x)\chi_{A}(x)\frac{u(x)}{\abs{u(x)}}d\mathcal{L}^{n}(x) = \int_{E}\abs{u(x)}\chi_{A \cap E_{h}}(x)d\mathcal{L}^{n}(x) \\
		& \geq C\mathcal{L}^n(A \cap E_h).
	\end{split}
	\]
	D'altra parte,
	\[
	\begin{split}
		\langle \varphi, \chi_{h}\chi_{A}\frac{u}{\abs{u}}\rangle = \langle \varphi, \chi_{A \cap E_{h}}\frac{u}{\abs{u}}\rangle \le \norm{\varphi}_{(L^{1}(E))'}\norm*{\chi_{A \cap E_{h}}\frac{u}{\abs{u}}}_{1} = \norm{\varphi}_{(L^{1}(E))'}\mathcal{L}^{n}\left( A \cap E_{h}\right),
	\end{split}
	\]
	così che
	\[ C\mathcal{L}^{n}(A\cap E_{n}) \le \int_{E}\abs{u(x)}\chi_{A \cap E_{h}}d\mathcal{L}^{n}(x) \le \norm{\varphi}_{(L^{1}(E))'}\mathcal{L}^{n}\left( A \cap E_{h}\right).  \]
	Dato che $C > \norm{\varphi}_{(L^{1}(E))'}$, deduciamo che $\mathcal{L}^{n}\left( A \cap E_{h}\right) = 0$. Dall'arbitrarietà di $h$, risulta che $\mathcal{L}^{n}\left( A \right) = 0$. Quindi $\abs{u(x)} < C$ q.o. in $E$, ovvero $u \in L^{\infty}(E) $.\\
	
	Mostriamo infine che per ogni $f \in L^{1}(E)$
	\begin{equation}\label{eq:representation_L1Linfty}
		\langle \varphi, f \rangle = \int_{E}f(x)u(x)d\mathcal{L}^{n}(x).
	\end{equation}
	Per fare ciò consideriamo l'operatore $T : L^1(E) \to L^\infty(E)$ definito da
	\[ T_{h}f(x) = \begin{cases}
		-h & \textup{se } f(x) \le -h, \\
		f(x) & \textup{se } \abs{f(x)} \le h, \\
		h & \textup{se } f(x) \ge h,
	\end{cases} \]
	e comunemente detto \emph{troncatura\index[analitico]{troncatura} $h$-esima di $f$}. Sostituendo nella \eqref{condizioneimportante} otteniamo
	\begin{equation}\label{dalimitizzare}
		\langle \varphi,\chi_{h} T_{h}f \rangle = \int_{E}\chi_{h}(x)T_{h}f(x)u(x)d\mathcal{L}^{n}(x).
	\end{equation}
	A questo punto, per $h \to +\infty$, $\chi_{h}(x)T_{h}f(x) \to f(x)$ q.o. in $E$ e
	\[ \abs{\chi_{h}(x)T_{h}f(x)}\le \abs{T_{h}f(x)} \le \abs{f(x)} \in L^{1}(E), \]
	così che per il Teorema della convergenza dominata $\chi_{h} T_{h}f \to f$ in $L^{1}(E)$. Possiamo quindi passare al limite nella \eqref{dalimitizzare} e ottenere \eqref{eq:representation_L1Linfty}.
	
	Supponiamo ora esista un'altra $\tilde{u} \in L^\infty(E)$ tale che per ogni $f \in L^{1}(E)$
	\[ \langle \varphi, f \rangle = \int_{E}f(x)\tilde{u}(x)d\mathcal{L}^{n}(x). \]
	Allora per ogni $f \in L^{1}(E)$ vale
	\[ \int_{E}f(x)\tilde{u}(x)d\mathcal{L}^{n}(x) = \int_{E}f(x)u(x)d\mathcal{L}^{n}(x),\]
	ovvero
	\[ \int_{E}(u(x)- \tilde{u}(x))f(x)d\mathcal{L}^{n}(x) =0.\]
	In particolare, la condizione sopra scritta vale per ogni $f \in C^{\infty}_{c}(E)$, così che $u - \tilde{u} = 0 $ q.o in $E$, da cui l'unicità.
\end{proof}

\begin{corollario}
	Siano $E \subseteq \R^{n}$ misurabile, $\varphi : L^{1}(E) \to \R$ lineare e continua e sia $u \in L^{\infty}(E)$ la rappresentazione di $\varphi$ nel senso del Teorema \eqref{dualediL1}. Allora 
	\[ 	\norm{\varphi}_{(L^{1}(E))'} = \norm{u}_{\infty}. \]
\end{corollario}

\begin{proof}
	Innanzitutto, per ogni $f \in L^{1}(E)$
	\[ \abs{\langle \varphi, f\rangle} = \int_{E}f(x)u(x)d\mathcal{L}^{n}(x) \le \norm{u}_{\infty}\norm{f}_{1} \]
	da cui $\norm{\varphi}_{(L^{1}(E))'}\le \norm{u}_{\infty} $. 
	
	Viceversa, dato $C > \norm{\varphi}_{(L^{1}(E))'} $, con passaggi analoghi a quelli del Teorema \eqref{dualediL1}, possiamo affermare che $\norm{u}_{\infty} \le C$. Scelto quindi $C = \norm{\varphi}_{(L^{1}(E))'} + \frac{1}{h}$, con $h \in \N\setminus \left\lbrace 0\right\rbrace $, passando al limite per $h \to +\infty$ otteniamo $\norm{u}_{\infty} \le \norm{\varphi}_{(L^{1}(E))'}$, da cui la tesi.
\end{proof}

\begin{proposizione}
	Sia $E \subseteq \R^{n} $ misurabile. Allora $L^{1}(E) \hookrightarrow (L^{\infty}(E))'$.
\end{proposizione}

\begin{proof}
	Sia $f \in L^{1}(E)$. Consideriamo l'applicazione lineare  $\psi : L^{\infty}(E) \to \R$ definita come
	\[ \psi(g) = \int_{E}^{}f(x)g(x)d\mathcal{L}^{n}(x).\]
	Si vede facilmente che
	\[ \abs{\psi(g)} \le \norm{f}_{1}\norm{g}_{\infty},\]
	così che $\psi$ è continua e di conseguenza $\psi \in (L^{\infty}(E))'$. L'identificazione è quindi fornita dall'applicazione $f \mapsto \psi$. Tuttavia, essa non è suriettiva.
	
	Infatti, si consideri $\xi \in E$ e si definisca l'applicazione $T : (C_{b}(E), \norm{\cdot}_{\infty}) \to \R$ come
	\[ \langle T, f \rangle = f(\xi). \]
	Essa è ben definita, lineare e tale che
	\[ \abs{\langle T, f \rangle} = \abs{f(\xi)} \le \max_E \abs{f} = \norm{f}_{\infty}, \]
	quindi continua. Siccome $C_{b}(E) \le L^{\infty}(E)$, per il Teorema di Hahn -- Banach (forma analitica) esiste $\mathfrak{T}: L^{\infty}(E) \to \R$ lineare e continua tale che 
	\[ \mathfrak{T}|_{C_{b}(E)} = T. \]
	Evidentemente $\mathfrak{T} \in (L^{\infty}(E))'$. Se, per assurdo, esistesse $\tilde{g} \in L^{1}(E)$ tale che $\mathfrak{T} = \psi(\tilde{g})$, allora si avrebbe
	\[ \Braket{\mathfrak{T}, f} = \int_{E}^{}f(x)\tilde{g}(x)d\mathcal{L}^{n}(x) \]
	per ogni $f \in L^{\infty}(E)$. Se $f \in C^{\infty}_{c}(E\setminus \left\lbrace \xi\right\rbrace )$, allora $f \in C_{b}(E)$ e $f(\xi) = 0$; ne segue che
	\[ \int_{E} f(x)\tilde{g}(x)d\mathcal{L}^{n}(x) = \langle T,f \rangle = f(\xi) = 0, \]
	da cui $\tilde{g}(x) = 0$ per q.o. $x \in E\setminus\left\lbrace \xi\right\rbrace $, ovvero $\tilde{g}(x) = 0$ per q.o. $x \in E$. Allora $\mathfrak{T} = 0$, ma $\mathfrak{T}|_{C_{b}(E)} = T \ne 0$, da cui l'assurdo. Pertanto l'inclusione è stretta, ovvero 
	\[ \psi(L^{1}(E)) \subsetneq (L^{\infty}(E))'. \qedhere\]
\end{proof}

\section{Convergenza debole in spazi di Banach}
\begin{definizione}\index[analitico]{convergenza! debole}
Siano $X$ uno spazio di Banach su $\K$ e $(x_h)$ una successione in $X$. Diciamo che \emph{$x_h$ converge debolmente a $x \in X$}, e lo indichiamo con $x_{h} \rightharpoonup x$, se per ogni $\varphi \in X'$
\[ \lim_{h} \braket{\varphi, x_h} = \braket{\varphi, x} \]
\end{definizione}

\begin{proposizione}
Siano $X$ uno spazio di Banach su $\K$ e $(x_h)$ una successione in $X$ ed $x \in X$ tale che $x_h \rightharpoonup x$ in $X$. Allora
\[ \norm{x} \le \underset{h}{\textup{lim inf}} \norm{x_h}. \]
\end{proposizione}
\begin{proof}
Innanzitutto, 
\[ \norm{x} = \sup \Set{\abs{\braket{\varphi,x}} :  \varphi \in X', \norm{\varphi}\le 1}. \]
Inoltre,
\[ \abs{\braket{\varphi,x_h}} \le \norm{\varphi}\norm{x_h} \]
e passando al lim inf, sfruttando la continuità
\[ \abs{\braket{\varphi, x}} = \underset{h}{\textup{lim inf}}\abs{\braket{\varphi,x_h}} \le \underset{h}{\textup{lim inf}}\norm{\varphi}\norm{x_h} \]
e passando al sup per $\norm{\varphi} \le 1$ otteniamo la tesi.
\end{proof}

\begin{proposizione}
Siano $X$ uno spazio di Banach su $\K$ e $(x_h)$ una successione in $X$. Se esiste $x \in X$ tale che $x_h \rightharpoonup x$ in X, allora $x$ è unico.
\end{proposizione}
\begin{proof}
	Omettiamo la dimostrazione.
\end{proof}
\begin{proposizione}
Siano $X$ uno spazio di Banach su $\K$ e $(x_h)$ una successione in $X$. Se esiste $x \in X$ tale che $x_h \to x$ in $X$, allora $x_h \rightharpoonup x$ in $X$.
\end{proposizione}
\begin{proof}
Omettiamo la dimostrazione.
\end{proof}
\begin{proposizione}
Siano $X$ uno spazio di Banach su $\K$ e $(x_h)$ una successione in $X$ limitata. Se
$x_h$ converge debolmente a qualche $x \in X$, allora $(x_{h})$ è limitata.
\end{proposizione}
\begin{proof}
	Omettiamo la dimostrazione.
\end{proof}
Il seguente è un analogo del Teorema di Bolzano--Weierstrass.
\begin{teorema}\emph{\textbf{(di Eberlein--Smulian)}}\index[analitico]{teorema!di Eberlein--Smulian}
	Sia $X$ uno spazio di Banach su $\K$. $X$ è riflessivo se e solo se ogni successione $(x_h)$ in X limitata ammette una sottosuccessione debolmente convergente.
\end{teorema}
\begin{proof}
	Omettiamo la dimostrazione.
\end{proof}
\begin{definizione}\index[analitico]{convergenza!debole*}
	Siano $X$ uno spazio di Banach su $\K$ e $(\varphi_h)$ una successione in $X'$. Diciamo che $\varphi_h$ \emph{converge debolmente*}, e lo indichiamo con $\varphi_h \overset{\ast}{\rightharpoonup} \varphi$\index[simboli]{$\overset{\ast}{\rightharpoonup}$} se per ogni $x \in X$
	\[ \braket{\varphi_h,x} \to \braket{\varphi, x}.\footnote{Praticamente è una convergenza puntuale.} \]
\end{definizione}
\begin{proposizione}\label{separabilitàe*}
	Siano $X$ uno spazio normato separabile su $\K$ e $(\varphi_h)$ una successione in $X'$ limitata. Allora esistono $\varphi \in X'$ ed una sottosuccessione $(\varphi_{h_{k}})$ tali che $\varphi_{h_{k}}\overset{\ast}{\rightharpoonup} \varphi $.
\end{proposizione}
\begin{proof}
	Omettiamo la dimostrazione.
\end{proof}

Veniamo al notevole caso degli spazi di Lebesgue.
\begin{osservazione}
Se $X = L^p(U)$, con $1 \le p < \infty$. Dal fatto che $(L^p(U))'\cong L^{p'}(U)$, per ogni $\varphi \in (L^p(U))'$, esiste $g \in L^{p'}(U)$ tale che 
\[ \braket{\varphi, f} = \int_{U}fgd\mathcal{L}^n. \]
Allora $f_h \rightharpoonup f$ se e solo se per ogni $g \in L^{p'}(U)$
\[ \lim_{h} \int_{U} f_hgd\mathcal{L}^n = \int_{U} fgd\mathcal{L}^n. \]
\end{osservazione}


\begin{proposizione}
Sia $p \in \left]   1,\infty\right[ $. Sia $(f_h)$ una successione in $L^p(U)$ tale che $f_h \to f$ \emph{q.o.} in $U$ e $f_h \rightharpoonup g$ in $L^p(U)$. Allora $g=f$ \emph{q.o.} in $U$.
\end{proposizione}
\begin{proof}
Innanzitutto, siccome $f_h \rightharpoonup g$, risulta che $(f_h)$ è limitata e per ogni $u \in C^\infty_{c}(U)$	
\[ \int_{U}f_h u d\mathcal{L}^n \to \int_U g u d\mathcal{L}^n. \]
Ora, siccome di fatto integriamo solo sul supporto di $u$, per la Proposizione \eqref{convergenzaLq}, per ogni $q < p$, $f_h u \to fu$ in $L^q(U)$. Siccome la convergenza forte implica quella debole, $f_hu \rightharpoonup fu $ in $L^q(U)$, e quindi, a meno di rinominazioni, per ogni $u \in C^\infty_{c}(U)$	
\[ \int_{U}f_h u d\mathcal{L}^n \to \int_U f u d\mathcal{L}^n. \]
A questo punto, per l'unicità del limite, deve essere per ogni $u \in C^\infty_{c}(U)$	
\[ \int_U g u d\mathcal{L}^n = \int_U f u d\mathcal{L}^n \]
ossia
\[ \int_U (g-f) u d\mathcal{L}^n = 0 \]
e quindi, per il Lemma di Du Bois-Reymond, $g = f$ q.o. in $U$.
\end{proof}

\begin{osservazione}\label{convergenzadeboleLp}
Consideriamo una successione $(u_h)$ in $L^p(U)$ limitata.
\begin{itemize}
	\item se $1 < p < \infty$, $L^p(U)$ è riflessivo quindi per il Teorema di Eberlein--Smulian esistono $u \in L^p(U)$ e una sottosuccessione $u_{h_{k}}$ tali che per ogni $\varphi \in L^{p'}(U)$
	\[ \int_{U}u_{h_{k}}\varphi d\mathcal{L}^n \to \int_{U}u\varphi d\mathcal{L}^n \]
	\item se $p = \infty$, siccome $L^\infty(U) \cong (L^1(U))'$, consideriamo $T: L^\infty(U) \to (L^1(U))'$ tale che 
	\[ \braket{Tu_h,\varphi} = \int_{U}u_h\varphi d\mathcal{L}^n. \]
	Essendo $L^1(U)$ separabile, dalla Proposizione \eqref{separabilitàe*} esistono  $G \in (L^1(U))'$ ed $Tu_{h_{k}}$ tali che per ogni $\varphi \in L^{1}(U)$
	\[ \braket{Tu_{h_{k}},\varphi} \to \braket{G,\varphi}\]
	ma deve essere $G = Tu$ per qualche $u \in L^\infty(U)$ e quindi  per ogni $\varphi \in L^{1}(U)$
	\[ \int_{U}u_{h_{k}}\varphi d\mathcal{L}^n \to \int_{U}u\varphi d\mathcal{L}^n \]
	\item se $p=1$, esiste $T: L^1(U) \to \left( C(\overline{U})\right)'$ tale che per ogni $g \in C(\overline{U})$,
	\[ \braket{Tu_h,g} = \int_U u_h g d\mathcal{L}^n. \]
	Essendo $C(\overline{U})$ separabile, per la Proposizione \eqref{separabilitàe*}, esiste $G \in \left( C(\overline{U})\right)'$ ed una sottosuccessione $(u_{h_k})$ tali che per ogni $g \in C(\overline{U})$ 
	\[ \int_U u_{h_k} g d\mathcal{L}^n \to \braket{G,g}. \]
	A questo punto, esiste una ed una sola misura di Radon $\nu$\footnote{Omettiamo la dimostriazione di questo fatto.} tale che
	\[ \braket{G,g} = \int_U gd\nu \]
	e quindi per ogni $g \in C(\overline{U})$
	\[ \int_U u_{h_k} g d\mathcal{L}^n \to \int_U gd\nu.\]
\end{itemize}
\end{osservazione}
\begin{proposizione}
	Sia $E \subseteq \R^{n} $ misurabile. Allora $L^{1}(E)$ non è riflessivo.
\end{proposizione}
\begin{proof}
	Consideriamo innanzitutto $(U_j)$ una successione decrescente di aperti di $E$ tale che $\mathcal{L}^n(U_j) \to 0$ e per ogni $j \in \N$
	\[ u_j = \frac{\chi_{U_j}}{\norm{\chi_{U_j}}_{L^1(E)}}. \]
	Chiaramente, $\norm{u_j}_{L^1(E)} = 1$, quindi $(u_j)$ è limitata in $L^1(E)$. Se, per assurdo, $L^1(E)$ fosse riflessivo, allora, per il Teorema di Eberlein--Smulian, esisterebbero $(u_{j_k})$ in $L^1(E)$ ed $u \in L^1(E)$ tali che $u_{j_k} \rightharpoonup u$ in $L^1(E)$. In altre parole, per ogni $\psi \in L^\infty(E)$
	\[ \int_E u_{j_k}\psi d\mathcal{L}^n \to \int_E u \psi d\mathcal{L}^n. \]
	Per ogni $h \in \N$ si ha $\chi_{U_h} \in L^\infty(E)$, quindi 
	\[ \int_E u_{j_k}\chi_{U_h} d\mathcal{L}^n \to \int_E u \chi_{U_h} d\mathcal{L}^n, \]
	ma $\chi_{U_{j_k}}\chi_{U_h} = \chi_{U_{j_k}\cap U_h}$, perciò, essendo anche $(U_{j_k})$ decrescente, per ogni $h \in \N$ esiste $\overline{k} \in \N$ tale che per ogni $k \ge \overline{k}$ si abbia $\chi_{U_{j_k}}\chi_{U_h} = \chi_{U_{j_k}}$. Allora
	\[  \int_E u_{j_k}\chi_{U_h} d\mathcal{L}^n \to 1 \]
	e, per unicità del limite,
	\[ \int_E u \chi_{U_h} d\mathcal{L}^n = 1. \]
	Ora, però, 
	\[ \lim\limits_{h} \norm{\chi_{U_h}}_{L^1(E)} = \lim\limits_{h} \mathcal{L}^n(U_h) = 0,\]
	quindi $\chi_{U_h} \to 0$ in $L^1(E)$. A meno di sottosuccessioni, per le proprietà degli spazi di Lebesgue, $\chi_{U_h} \to 0$ q.o. in $L^1(E)$. In particolare, $u\chi_{U_h} \to 0$ q.o. in $L^1(E)$ e $u\chi_{U_h} \le u \in L^1(E)$ per cui, applicando il Teorema della convergenza dominata,
	\[  \lim\limits_{h} \int_E u\chi_{U_h} d\mathcal{L}^n = 0, \]
	da cui l'assurdo.
\end{proof}

\section*{Riepilogo proprietà degli Spazi di Lebesgue}
Sia $E$ un sottoinsieme misurabile di $\R^{n}$. Possiamo così riepilogare le proprietà di $(L^{p}(E), \norm{\;}_{p})$.

\hfill

\begin{center}
	\begin{tabular}{lccc}
		
		& $p=1$ & $1<p<\infty$ & $p=\infty$ \\
		
		Spazio di Banach & \checked & \checked & \checked \\
		
		Spazio di Banach separabile & \checked & \checked &  \\
		
		Spazio di Banach riflessivo &  & \checked &  \\
		
	\end{tabular}
\end{center}

\section{Approssimazione polinomiale di funzioni continue}
Il seguente risultato è stato dimostrato da Karl Weierstrass nel 1855.
\begin{teorema}\emph{\textbf{(di approssimazione di Weierstrass)}}\index[analitico]{teorema!di approssimazione di Weierstrass}
Siano $a,b \in \R$, con $a<b$, ed $f : [a,b] \to \R$ un'applicazione continua. Allora esiste una successione di polinomi $(p_h)$ tale che $p_h \to f$ uniformemente su $[a,b]$.
\end{teorema}
\begin{proof}
Innanzitutto, a meno di una traslazione ed un riscalamento possiamo supporre $[a,b] = [0,1]$. Inoltre, senza perdita di generalità possiamo supporre che $f(0) = f(1) = 0$.\footnote{Consideriamo l'applicazione 
	\[ \tilde{f}(x) = f(x) - f(0) -x(f(1) - f(0)).  \]
	Chiaramente, $\tilde{f}$ realizza l'ipotesi richiesta. Se dimostriamo che la tesi vale per $\tilde{f}$, essendo quest'ultima $f$ privata di un polinomio, allora sarà vera anche per $f$. Non è restrittivo considerare quindi fin da subito $f$ con l'ipotesi menzionata.} Possiamo inoltre, a meno di un prolungamento continuo, pensare che $f = 0$ fuori da $[0,1]$. In particolare, $f$ è uniformemente continua.
	
	Consideriamo il polinomio $Q_n$ tale che
	\[ Q_n(x) = c_n(1-x^2)^n, \]
	dove $c_n \in \R$ tale che
	\[ \int_{-1}^1 Q_n d\mathcal{L}^1 = 1. \]
	In particolare, siccome $(1-x^2)^n \ge 1-nx^2$ su $[0,1]$,\footnote{Sia $u:[0,1]\to \R$ tale che
	\[ u(x) = (1-x^2)^n - 1+nx^2. \]
	Sicuramente $u$ è continua su $[0,1]$ e derivabile su $\left] 0,1\right[ $. Inoltre $u(0) = 0$ e 
	\[ u'(x) = -2nx (1-x^2)^{n-1} +2nx = 2nx(1-(1-x^2)^{n-1}) \ge 0, \]
	quindi $u(x) \ge 0$.}
	\begin{align*}
	\int_{-1}^1 (1-x^2)^n dx &= 2 \int_0^1 (1-x^2)^ndx \ge 2\int_{0}^{\frac{1}{\sqrt{n}}}(1-x^2)^n dx \ge 2\int_{0}^{\frac{1}{\sqrt{n}}}(1-nx^2) dx =  \\ 
	&=\frac{4}{3\sqrt{n}} > \frac{1}{\sqrt{n}},
	\end{align*}
	quindi $c_n < \sqrt{n}$. Inoltre, dato $\delta \in \left] 0,1\right[ $, su $[-1,-\delta]\cup [\delta,1]$
	\[ Q_n(x) < \sqrt{n} (1-\delta^2)^n, \]
	quindi $Q_n \to 0$ uniformemente su $[-1,-\delta]\cup [\delta,1]$.
	
	Consideriamo l'applicazione $p_n: [0,1]\to \R$ tale che
	\[ p_n(x) = \int_{-1}^1 f(x+t)Q_n(t)dt. \]
	Innanzitutto, siccome $f$ è nulla fuori da $[0,1]$,
	\[ p_n(x) = \int_{-x}^{1-x}f(x+t)Q_n(t) dt = \int_0^1 f(t) Q_n(t-x) dt, \]
	ossia $p_n$ è un polinomio in $x$. Sia $M = \max f$. Per ogni $\varepsilon>0$, sia $\delta >0$ tale che per ogni $x,y \in [0,1]$ con $\abs{x-y} <\delta$, si abbia $\abs{f(x)-f(y)}<\frac{\varepsilon}{2}$. Senza perdita di generalità possiamo supporre $\delta <1$. A questo punto, per ogni $x \in [0,1]$, se $t \in \left] -\delta,\delta\right[ $ si ha $\abs{x+t-x} = \abs{t}<\delta$, allora
	\begin{align*}
		\abs{p_n(x)-f(x)}&\le \int_{-1}^1 \abs{f(x+t)-f(x)}Q_n(t) dt = \int_{-1}^{-\delta}\abs{f(x+t)-f(x)}Q_n(t) dt + \\
		&+ \int_{-\delta}^{\delta}\abs{f(x+t)-f(x)}Q_n(t) dt+ \int_{\delta}^{1}\abs{f(x+t)-f(x)}Q_n(t) dt <\\
		&<2M \int_{-1}^{-\delta} Q_n d\mathcal{L}^1 + \frac{\varepsilon}{2}\int_{-\delta}^{\delta}Q_nd\mathcal{L}^1 +2M \int_{\delta}^{1} Q_n d\mathcal{L}^1 \le \\
		&\le 2M\left( \int_{-1}^{-\delta} Q_n d\mathcal{L}^1 + \int_{\delta}^{1} Q_n d\mathcal{L}^1\right)  + \frac{\varepsilon}{2}\int_{-1}^{1}Q_nd\mathcal{L}^1 = \\
		&=2M\left( \int_{-1}^{-\delta} Q_n d\mathcal{L}^1 + \int_{\delta}^{1} Q_n d\mathcal{L}^1\right)  + \frac{\varepsilon}{2}
	\end{align*}
	e, per il fatto che $Q_n \to 0$ uniformemente su $[-1,-\delta]\cup [\delta,1]$, esiste $N >0$ tale che per ogni $n \ge N$ si abbia $Q_n\le \frac{\varepsilon}{8M}$ su $[-1,-\delta]\cup [\delta,1]$, da cui 
	\[ 	\abs{p_n(x)-f(x)}< \varepsilon, \]
	quindi $p_n \to f$ uniformemente su $[0,1]$, da cui la tesi.
\end{proof}

\begin{esempio}
Consideriamo $f: [0,1]\to \R$ continua tale che per ogni $n \in \N$
\[ \int_{0}^{1}f(x)x^n dx = 0. \]
Mostriamo che $f = 0$.

Innanzitutto, per linearità, per ogni polinomio $p$
\[ \int_{0}^{1}f p d\mathcal{L}^1 = 0. \]
Ora, per il Teorema di approssimazione di Weierstrass, esiste $(p_h)$ tale che $p_h \to f$ uniformemente su $[0,1]$, quindi 
\[ \int_{0}^{1} f^2 d\mathcal{L}^1 = \lim\limits_{h} \int_{0}^{1} fp_h d\mathcal{L}^1 = \lim\limits_{h} 0 = 0, \]
da cui, essendo $f$ continua, $f=0$.
\end{esempio}

Miriamo ora a presentare la generalizzazione del risultato di Weierstrass proposta da Mashall Stone nel 1937: a malincuore dobbiamo parlare di algebre.
\begin{definizione}
Consideriamo $(X,\norm{\;})$ uno spazio normato su $\R$ ed $E \subseteq X$. Diciamo che $\mathcal{A} \subseteq C_b(E; \R^n)$ è un'\emph{algebra}\index[analitico]{algebra} (su $\R$) se valgono i seguenti fatti:
\begin{enumerate}[$(a)$]
\item se $f,g \in \mathcal{A}$, allora $f+g \in \mathcal{A}$,
\item se $f,g \in \mathcal{A}$, allora $fg \in \mathcal{A}$,
\item se $k \in \R$ e $f \in \mathcal{A}$, allora $kf \in \mathcal{A}$.
\end{enumerate}
\end{definizione}

\begin{lemma}
Consideriamo $(X,\norm{\;})$ uno spazio normato su $\R$ ed $E \subseteq X$. Se $\mathcal{A} \subseteq C_b(E; \R^n)$ è un'algebra, allora $\overline{\mathcal{A}}$ è un'algebra.
\end{lemma}
\begin{proof}
Siano $f,g \in \overline{\mathcal{A}}$ e $k \in \R$. Esistono $(f_h),(g_h)$ in $\mathcal{A} \subseteq C_b(E; \R^n)$ tali che $f_h \to f$, $g_h \to g$ uniformemente su $E$. Essendo $C_b(E;\R^n)$ chiuso in $\mathcal{B}(E;\R^n)$, ne deduciamo che $f,g \in C_b(E; \R^n)$. A questo punto, chiaramente $f_h + g_h \to f+g$ e $kf_h \to kf$ uniformemente su $E$, quindi $f+g, kf \in \overline{\mathcal{A}}$. Inoltre, essendo tutte le applicazioni limitate, anche $f_hg_h \to fg$ uniformemente su $E$, da cui $fg \in \overline{\mathcal{A}}$.
\end{proof}

\begin{definizione}
Consideriamo $(X,\norm{\;})$ uno spazio normato su $\R$ ed $E \subseteq X$. Diciamo che un'algebra $\mathcal{A} \subseteq C_b(E; \R^n)$ \emph{separa i punti}\index[analitico]{algebra!che separa i punti} se per ogni $x,y \in E$ esiste $f \in \mathcal{A}$ tale che $f(x) \ne f(y)$.
\end{definizione}

\begin{definizione}
Consideriamo $(X,\norm{\;})$ uno spazio normato su $\R$ ed $E \subseteq X$. Diciamo che un'algebra $\mathcal{A} \subseteq C_b(E; \R^n)$ \emph{non si annulla}\index[analitico]{algebra!che non si annulla} su $E$ se per ogni $x \in E$ esiste $f \in \mathcal{A}$ tale che $f(x) \ne 0$.
\end{definizione}

\begin{lemma}\label{tecnicoperSW}
Consideriamo $(X,\norm{\;})$ uno spazio normato su $\R$ ed $E \subseteq X$. Se $\mathcal{A} \subseteq C_b(E; \R)$ è un'algebra che separa i punti e non si annulla su $E$, allora per ogni $x,y \in E$ e per ogni $c_1,c_2 \in \R$ esiste $f \in \mathcal{A}$ tale che $f(x) = c_1$ e $f(y) = c_2$.
\end{lemma}
\begin{proof}
Dal fatto che $\mathcal{A}$ separa i punti, esiste $g \in \mathcal{A}$ tale che $g(x) \ne g(y)$. Inoltre, dal fatto che $\mathcal{A}$ non si annulla su $E$, esistono $h,k \in \mathcal{A}$ tale che $h(x),k(y) \ne 0$. Posto
\begin{align*}
 u &= gk -g(x)k, & v= gh-g(y)h,
\end{align*}
basta considerare 
\[ f = \frac{c_2u}{u(y)} + \frac{c_1v}{v(x)}.\qedhere \]
\end{proof}

\begin{teorema}\emph{\textbf{(di Stone--Weierstrass)}}\index[analitico]{teorema!di Stone--Weierstrass}
Consideriamo $K \subseteq \R^n$ compatto. Se $\mathcal{A} \subseteq C(K; \R^m)$ è un'algebra che separa i punti e non si annulla su $K$, allora $\overline{\mathcal{A}} = C(K; \R^m)$.
\end{teorema}
\begin{proof}
Innanzitutto, a meno di ragionare per componenti, è sufficiente considerare il caso $\mathcal{A} \subseteq C(K; \R)$. Sicuramente, $\overline{\mathcal{A}} \subseteq C(K; \R^m)$. 

Mostriamo, prima di tutto, che se $f \in \overline{\mathcal{A}}$, allora $\abs{f} \in \overline{\mathcal{A}}$. In primo luogo, per il Teorema di Weierstrass, $f$ è limitata, ossia esiste $a \in \R$ tale che per ogni $x \in K$ si abbia $\abs{f(x)}\le a$. Per il Teorema di approssimazione di Weierstrass, esiste $(p_n)$ convergente uniformemente a $\abs{x}$ su $[-a,a]$: in altre parole, per ogni $\varepsilon > 0$ esistono $c_1,\dots,c_{n_\varepsilon} \in \R$ tali che per ogni $x \in [-a,a]$
\[ \abs*{\sum_{i=1}^{n_\varepsilon}c_ix^i- \abs{x}}<\varepsilon.\footnote{Siccome l'applicazione valore assoluto si annulla in zero, la successione di polinomi approssimanti può essere presa senza termini noti.} \]
Siccome $g_\varepsilon = \sum_{i=1}^{n_\varepsilon}c_if^i \in \overline{\mathcal{A}}$ e $\abs{g_\varepsilon-\abs{f}}<\varepsilon$ su $K$, ne deduciamo che $\abs{f} \in \overline{\mathcal{A}}$. In particolare, date $f,g \in \overline{\mathcal{A}}$, 
\begin{align*}
	\max\left\lbrace f,g\right\rbrace  &= \frac{1}{2}(f+g + \abs{f-g})\in \overline{\mathcal{A}}, & \min\left\lbrace f,g\right\rbrace  &= \frac{1}{2}(f+g - \abs{f-g})\in \overline{\mathcal{A}}
\end{align*}
e, per induzione, lo stesso vale per il massimo ed il minimo tra un numero finito di elementi di $\overline{\mathcal{A}}$.

Fissiamo $f \in C(K;\R)$ ed $\varepsilon>0$. Mostriamo che per ogni $x \in K$ esiste $g_x \in \overline{\mathcal{A}}$ tale che $g_x(x) = f(x)$ e per ogni $t \in K$ si abbia $g_x(t)\ge f(t)-\frac{\varepsilon}{2}$. Per il Lemma \eqref{tecnicoperSW}, per ogni $y \in K$ esiste $h_y \in \mathcal{A}\subseteq \overline{\mathcal{A}}$ tale che $h_y(x) = f(x)$ e $h_y(y) = f(y)$. Siccome $h_y - f$ è continua, per il Teorema di di permanenza del segno, esiste $U_y$ intorno aperto di $y$ tale che $h_y \ge f -\frac{\varepsilon}{2}$ su $U_y$. Siccome 
\[ K \subseteq \bigcup_{y \in K} U_y, \]
per compattezza, esistono $y_1,\dots,y_m$ tali che 
\[ K \subseteq \bigcup_{j=1}^m U_{y_j}. \]
Consideriamo $g_x = \max \left\lbrace h_{y_1},\dots, h_{y_m}\right\rbrace \in \overline{\mathcal{A}}$. Allora, 
\[ g_x(x) = \max \left\lbrace h_{y_1}(x),\dots, h_{y_m}(x)\right\rbrace = f(x) \]
e per ogni $t \in K$ esiste $i =1,\dots, m$ tale che $t \in U_{y_i}$, quindi 
\[ g_x(t) \ge h_{y_i}(t) \ge f(t)-\frac{\varepsilon}{2}, \]
quindi $g_x$ soddisfa i requisti richiesti.

A questo punto, sempre per il Teorema di permanenza del segno, per ogni $x \in K$ esiste un intorno aperto $V_x$ di $x$ tale che $g_x \le f+\frac{\varepsilon}{2}$ su $V_x$. Siccome 
\[ K \subseteq \bigcup_{x \in K} V_x, \]
per compattezza, esistono $x_1,\dots,x_N$ tali che 
\[ K \subseteq \bigcup_{j=1}^N V_{x_j}. \]
Preso $h = \min \left\lbrace g_{x_1},\dots g_{x_N}\right\rbrace \in \overline{\mathcal{A}}$, si ottiene che $\abs{h - f} \le \frac{\varepsilon}{2} <\varepsilon$, da cui la tesi.
\end{proof}

\begin{osservazione}
A ben vedere, non serve conoscere il Teorema di approssimazione di Weierstrass per dimostrare il risultato generale. Nella dimostrazione utilizziamo solo il fatto che l'approssimazione polinomiale valga per l'applicazione valore assoluto: in particolare, è sufficiente sapere che, per ogni $a >0$, la successione $(p_n)$ definita ricorsivamente da 
\begin{align*}
		p_0 &= 0, & p_{n+1}(x) &= p_n(x) + \frac{1}{2}(x^2 -p_n(x)^2),
\end{align*}
converge uniformemente a $\abs{x}$ su $[-a,a]$.
\end{osservazione}

\begin{osservazione}
Consideriamo $E \subseteq \R^n$ compatto. L'algebra delle applicazioni a componenti polinomiali separa i punti\footnote{Basta considerare l'applicazione identica, che ha componenti polinomiali.} e non si annulla su $E$\footnote{Basta considerare un'applicazione costante con almeno una componente non nulla.}, quindi per ogni $f \in C(E; \R^n)$ esiste una successione $(p_h)$ di applicazioni a componenti polinomiali tale che $p_h \to f$ uniformemente su $E$.
\end{osservazione}

\section{Spazi di Lebesgue che coinvolgono il tempo}
Prima di tutto necessitiamo di generalizzare la teoria dell'integrazione di Lebesgue.
\begin{definizione}
Consideriamo $(X,\norm{\;})$ uno spazio di Banach su $\R$ e $T>0$. Diciamo che un'applicazione $s: [0,T] \to X$ è \emph{semplice}\index[analitico]{applicazione!semplice} se
\[ s(t) = \sum_{i=1}^{m}\chi_{E_i}(t)u_i,  \]
dove $E_1,\dots,E_m \subseteq [0,T]$ misurabili e $u_1,\dots,u_m \in X$. 
\end{definizione}

\begin{definizione}
Consideriamo $(X,\norm{\;})$ uno spazio di Banach su $\R$ e $T>0$. Diciamo che un'applicazione $f: [0,T] \to X$ è \emph{fortemente misurabile}\index[analitico]{applicazione!misurabile!fortemente} se esiste una successione $(s_h)$ di funzioni semplici tali che $s_h \to f$ \emph{q.o.} in $[0,T]$.
\end{definizione}

\begin{definizione}
Consideriamo $(X,\norm{\;})$ uno spazio di Banach su $\R$ e $T>0$. Se $s: [0,T] \to X$ è un'applicazione semplice, definiamo 
\[ \int_0^T sd\mathcal{L}^1 = \sum_{i=1}^{m} \mathcal{L}^1(E_i)u_i. \]
\end{definizione}

\begin{definizione}
Consideriamo $(X,\norm{\;})$ uno spazio di Banach su $\R$ e $T>0$. Diciamo che un'applicazione $f: [0,T] \to X$ fortemente misurabile è \emph{sommabile}\index[analitico]{applicazione!sommabile} se esiste una successione $(s_h)$ di funzioni semplici tali che 
\[ \lim\limits_{h} \int_{0}^{T} \norm{s_h(t)-f(t)}d\mathcal{L}^1(t) = 0.  \]
\end{definizione}

\begin{definizione}
	Consideriamo $(X,\norm{\;})$ uno spazio normato su $\R$ e $T>0$. Se $f: [0,T] \to X$ è un'applicazione sommabile, definiamo 
	\[ \int_0^T fd\mathcal{L}^1 = \lim\limits_{h} \int_0^T s_hd\mathcal{L}^1 \]
\end{definizione}

\begin{definizione}
Consideriamo $(X,\norm{\;})$ uno spazio di Banach su $\R$, $p \in [1,\infty]$ e $T>0$. Chiamiamo $L^p(0,T;X)$\index[simboli]{$L^p(0,T;X)$} l'insieme delle (classi di equivalenza) di applicazioni fortemente misurabili $u: [0,T]\to X$ tali che, se $p \in \left[ 1,\infty\right[$, 
\[ \int_0^T \norm{u(t)}^p d\mathcal{L}^1 <\infty \]
oppure, nel caso $p=\infty$,
\[ \underset{t \in [0,T]}{\textup{ess sup}}\; \norm{u(t)} <+\infty. \]
\end{definizione}

\begin{definizione}
Consideriamo $(X,\norm{\;})$ uno spazio di Banach su $\R$, $p \in [1,\infty]$ e $T>0$. Data $u \in L^p(0,T;X)$, poniamo 
	\[ \norm{u}_{L^p(0,T;X)} = \begin{cases}
		\left( \displaystyle \int_0^T \norm{u(t)}^p d\mathcal{L}^1\right) ^{\frac{1}{p}} & 1 \le p < \infty, \\
		\underset{t \in [0,T]}{\textup{ess sup}}\; \norm{u(t)} & p=\infty.
	\end{cases} \]
	
\end{definizione}